\subsection{Charge carrier velocity\label{LadungGeschw}}
\Aufgabe{
    A current $I = 8 \ A$ flows through a copper wire with a cross-sectional area $A = 1 \ mm^2$ and a length $L = 10 \ m$. One $mm^3$ of copper contains $8.5\cdot10^{19}$ atoms.
    It can be assumed that 1 electron per atom is involved in charge transport. $(\vartheta = 20 \ ^\circ C)$.
    \newline
    \begin{enumerate}[label=\alph*)]
        \item Determine the drift velocity of the electrons in the copper wire.
        \item What is the electric field strength $E$ in the copper wire?
        \item How high is the voltage drop $U$ in this copper wire?
        \item What is the resistance of the copper wire under the specified boundary conditions? What resistance does the wire assume when heated to $\vartheta_{\mathrm{w}} = 180 \ ^\circ C$?
        How much is the percentage increase in resistance?
    \end{enumerate}
}

\Loesung{
	Hier entsteht eine Musterlösung...
}